\section{Realizing the expansion on a fixed network}\label{sec:assembly}

Different summands in Section~\ref{sec:full} can finish with different
wires and different tables. To represent their sum on one network, we
encode a choice of summand, together with an assignment to its final
wires, as a transcript. A local check determines which transcripts
contribute, and a dictionary supplies their tensor entries. The size
bound depends on how the host is shared. All blocks at one host
position use the same address. The bits carried on a wire vary
independently across blocks.

\subsection{From a transcript to a tensor}

Let $K$ be the working extension of $k$, and put
$h=\lceil\log_2 n\rceil$. Fix $P\le n^{O(1)}$ numbered processors.
Each has binary advice of at most $h^c$ bits and a fixed ordered list
of exterior ports. Its local record includes guesses and at most
$h^c$ incoming and outgoing messages. A message consists of an
$O(h)$-bit key identifying source, destination and port, together with
a bit string. Communication is valid if every used key occurs exactly
once at each end with the same string.

One fixed finite bit-RAM program checks the local record against the
advice and exterior bits. It may emit strings of length at most $U$ as
arguments to a common dictionary $W$ with values in $K$. Its
decisions depend only on the strings. For fixed exterior bits, the
transcript tensor is the sum over accepted transcripts of the product
of all values returned by $W$. The record format, including length
fields and padding, is fixed. All size-dependent program data enter
through the advice.

The charge at a processor includes its deterministic running time,
advice, exterior bits, guesses, messages with headers, and emitted
strings. Count at least one unit per record or item. This charge is
determined by the record. A local address or time uses
$(\log h)^{O(1)}$ bits. Bit arithmetic on words of that length adds at
most a factor $(\log h)^{O(1)}$ to the bounds below.

\begin{lemma}[Realizing a transcript]\label{lem:fixed-transcript}
Suppose the charge of each local record is at most
$V_{\max}=2^{(\log h)^{O(1)}}$. Let $V\le n^{O(1)}$ bound the total
charge of every accepted transcript and the volume of the fixed
advice and exterior interface, including one unit per processor. If
$U\le h$ and the wide-record bound of Theorem~\ref{thm:sort} holds,
then for every fixed $\tau>0$ the transcript tensor has a fixed
network of size at most
\[
                          h^\tau V+P h^{C_0}.
\]
The gates are trivalent or are dictionary tables of arity
$O(1+U+\log h)$. Their labels use $(\log h)^{O(1)}$ bits. The
constant $C_0$ depends only on $\tau,c$, the polynomial bounds in $n$
and the wide-record estimate. The exponent in $V_{\max}$ occurs only
in factors $(\log h)^{O(1)}$ before absorption.
\end{lemma}

\begin{proof}
\emph{Allocate space by record length.}
For each dyadic capacity $t$, from $1$ to the least power of two
at least $\min(V,V_{\max})$, reserve
$N_t=\min(P,\lfloor 2V/t\rfloor)$ slots. A processor with charge $s$
uses the least $t\ge s$. Every record of this class has charge
greater than $t/2$, so $N_t$ slots suffice. Guess their processor
identifiers and records. Check the identifier range and the charge,
require occupied slots to form an increasing prefix in identifier
order, and pin all unused entries and padding. The two relevant
capacity bounds are
\begin{equation}\label{eq:asm-seats}
 \sum_t tN_t=O\bigl(V(1+\log V_{\max})\bigr),\qquad
 \sum_t N_t=O\bigl(P(1+\log V_{\max})\bigr).
\end{equation}

\emph{Check the computation and its dictionary arguments.}
In a slot of capacity $t$, guess the deterministic execution for at
most $t$ steps. Time-ordered checks verify its instructions, registers,
counter, addresses and halt flag, with fixed padding after halt. To
check the memory trace, sort the reads, writes and initial contents by
address and execution order. Scan each address to verify that a read gives the
most recent value, starting at zero for unused cells. This is the
trace construction of \citet[Section~\ext{G}{app:trace}]{oaG}.
A comparison network on $O(t)$ accesses uses $O(t\log^2(2t))$
comparisons of $(\log h)^{O(1)}$-bit words. Long inputs, advice,
messages and emitted strings occupy addressed arrays and are
processed in pieces. The local cost is $t(\log h)^{O(1)}$ plus
$h^{O(1)}$ per slot for long identifiers and framing.

Within the slot, group the emitted strings by least dyadic padded
length $u$, counting an empty string as length one. Reserve
$\lfloor2t/u\rfloor$ calls in that class and take them in issuance
order. The checker writes each emitted bit at the address given by
its length class, ordinal and bit offset. The same memory-trace test
checks these arrays against the inputs of the installed calls,
including their active flags, actual lengths and zero padding. Each
address has $O(1+\log t+\log h)$ bits. Accounting for all the padded
seats adds $t(\log h)^{O(1)}$ to the checking cost. A seat returns
the value of $W$ on its indicated active string and returns $1$ when
inactive. Its table is named by the padded length. The arity is
$O(1+U+\log h)$ and the total table size per slot is
$O(t\log(2+U)\log(2+h))$.

\emph{Match the messages.}
Fix $w=\lceil h^{C_1}\rceil$, where $C_1>1$ will be chosen for
the wide-record sort. Split each message into $w$-bit chunks, padding
the last one and reserving one for an empty message. Its chunk record
contains an active flag, key, direction, total length, chunk number
and payload. Add a fixed sender for each processor. Its key belongs
to a reserved namespace and supplies that processor's advice and exterior bits with their
component lengths. An occupied slot must receive the whole supply
for its claimed identifier even when the supply is empty.

Check each outgoing source and incoming destination against the slot
identifier, and check the splitting, piece order, lengths and padding. Sort the
records by active flag, key, chunk number and direction. In the active
part, require every key and chunk to occur as one send and one receive
of equal length and payload. A duplicate message on one side collides
at chunk zero. In particular each private supply is matched to exactly
one slot, so every processor occurs once. For a given input the
stable-sort relation has just one output, which can be summed out after
these checks. Its zero-one tensor over $k$ extends to the same relation
over $K$.

A slot reserves $O(t/w+h^c+1)$ chunk records. The $O(h)$-bit header
is copied once per chunk, so constructing them costs
$(t+h^{O(1)})(\log h)^{O(1)}$. Charge a fixed sender for its actual
supply and one partial chunk. By Equation~\eqref{eq:asm-seats} the
reserved bit volume, including padding at message boundaries, is
\[
                     V(\log h)^{O(1)}+P h^{O(1)} .
\]
Pad the list to a power of two. Its length is polynomial in $n$, so
its logarithm is $O(h)$. Theorem~\ref{thm:sort} with tolerance
$\tau/3$ permits a fixed $C_1$ large enough that sorting costs at
most $O(h^{\tau/3})$ times this volume. The subsequent scan has
linear size. Equation~\eqref{eq:asm-seats} gives the same volume
bound for the local computations and the dictionary tables. Absorb
the powers of $\log h$ into $h^\tau$ and the per-processor costs
into $P h^{C_0}$.

The checking and sorting gates are trivalent over $k$, and each
dictionary gate is named by its length class. These names have the
stated length. For each accepted transcript the slot capacities,
identifiers, order, traces, calls and chunks are uniquely specified.
Conversely the local and communication checks recover an accepted
transcript. Contracting the installed relations therefore gives its
weighted sum exactly over $K$.
\end{proof}

\subsection{The transcript of the rounds}

We now apply Lemma~\ref{lem:fixed-transcript} to the expansion in
Section~\ref{sec:full}. Initially there are $S$ blocks of $T_0$
copies, so $d=T_0S$. At every stage of the $j$ rounds there are
$N_b\ge R$ blocks on a host with $H$ allocated positions. Give each
position one processor, including at the terminal stage, and give
each step $O(H)$ more processors for balanced aggregation trees. An
occurrence consists of one position in one block.

\paragraph{One address for all blocks.}
A position holds the ordered list of its host neighbors and an array
of occurrence records in block order. Its advice specifies the
parent and child naming conventions, integer parameters, field
presentation, and the initial convention for ordering exterior ports. At an
initial position it also gives the input table label. Exterior bits
enter through fixed ports in stack order.

In one step every host degree is bounded in terms of the fixed round
design. A ball of radius $O_D(1+\log r)$ thus contains $r^{O_D(1)}$
positions. Exchange neighbor lists to that radius, remove repetitions
by identifier and order the resulting names. This locates each
possible child and contact by an index in a local list, at a cost of
$h^{O(1)}$ common operations per position. A message to a contact
can now carry an entire array of occurrence records under one
$O(h)$-bit header. Within that array only local indices are needed.
There are $h^{O(1)}$ phases and declarations per position. By
Equation~\eqref{iter:labels}, handling the arrays costs
$N_b r^{O_D(1)}(\log h)^{O(1)}h^{O_D(a)}
=N_b h^{O_D(a)+o(1)}$ per position. The coefficients in $O_D(a)$
are independent of $a,b,\beta$.

For decisions shared by the host, each position computes its local
contributions to the candidate scores. Sum the integer arrays along
a balanced tree and broadcast the first minimizing candidate. At
height $i$ the width of a sum of initially $z$-bit entries is at
most $z+i$, while the number of positions falls by a factor of order
$2^i$. Since
$\sum_{i\ge0}2^{-i}(z+i)^{O(1)}=O((1+z)^{O(1)})$, aggregation
preserves the preceding bit-cost bound. At the root allow an
additional $2^{(\log h)^{O(1)}}$ operations to choose the
minimizer. A greedy coloring of a given host can be checked locally
as well. In fixed identifier order, require each color to be the
least one absent from its lower neighbors. Induction on identifiers
makes the checked assignment unique.

\paragraph{Checking the changes of host.}
Each of the operations in a round supplies its contacts from these
neighbor lists and its common parameters from additive scores.
\begin{enumerate}
\item For colony packing, try the multipliers of the $r$ slots,
propagate the overloaded marks to ancestors and neighbors, and sum
the candidate colony counts. A propagation moves upward by at least
one dyadic level, so the radius is $O_D(1+\log r)$. Relative
interval endpoints, loads and dyadically scaled charges all have
short encodings. Summing the local scores also chooses the affine
coloring for whole-item pairing.
\item At a local concentration step, the colony interval and its
old contacts give the possible neighbors of each sphere. In the
checked color order, the movement and interpolation choices of
Lemma~\ref{lem:local} specify the new tables, withdrawn scalars
and ordered bundles. Its description bounds pay for the choices and
labels. The port orders identify the exposed sublists.
\item In a wire-saving subround the source colony and bin determine
the seats of its matched units. Record one summand choice for each
actual gadget. Sum the surviving replacement widths under each of
the two contractions and take the same minimizer for all blocks.
The widths before identifying coincident endpoints suffice for this
decision. Bounded neighborhoods of the colonies then determine the
new contacts, port orders and table labels.
\item For a reset, compute each position's mean load over the blocks.
The smoothed mean has floor $FB/r$ and a propagated contribution
decays by at least a factor of two per edge. The starting load is at
most $B$, so a neighborhood of radius $O(1+\log r)$ determines the
bases and their capacities.

For the common mixer, let $x_i$ be the normalized load at a position
or the normalized width at an edge in block
$i\in\mathbb F_2^{\log_2 N_b}$. At prefix length $s$ the part of
the collision estimator of Section~\ref{subsec:reset} that depends
on a candidate matrix $A$ is a multiple of
\[
 \sum_{z\ne0}\boldsymbol{1}_{(Az)_{1:s}=0}
                     \sum_i x_i x_{i+z}.
\]
Two indices have the same output prefix exactly when their
difference $z$ satisfies the displayed condition. The diagonal and
centering terms are constant in $A$, so they can be omitted in
choosing a minimizer. The denominators of the remaining $x_i$ are
cleared by $FB$ times the fixed dyadic load denominator. The level
and overflow weights are integers after fixed scaling. In
particular the choice uses integer arithmetic in every
characteristic.

Compute the autocorrelations $\sum_i x_i x_{i+z}$ simultaneously by
squaring the Walsh coefficients and inverting the transform. This
takes $O(N_b\log N_b)$ operations on $(\log h)^{O(1)}$-bit local
integers. Aggregate the coefficients by difference and prefix
length. At the root one can enumerate all invertible matrices in
dimension $\log_2 N_b=O((\log h)^2)$ within
$2^{(\log h)^{O(1)}}$ bit operations. The first minimizer satisfies
the bound proved by averaging. Positions permute their block arrays
accordingly. The bounded length ratios of intersecting leaves and
the bounded number of units per leaf then allow the terminal
assignments, rotor settings and new labels to be determined from
neighboring records.
\item The final splitting trees for tethers have bounded size,
depending on $D$, per possible station and contact. Adjacent
positions determine the active sublists. An edge between new
children is supported by a bounded path between their parents. In
the next phase the children use their own host adjacency, so every
neighborhood collected is local to its current step.
\end{enumerate}

Declare messages in a fixed order of phases, ports and blocks.
Enforce each successive array and common host from the preceding
data and the local summand choices. Each term of a local identity,
including a term of weight zero, has one prescribed encoding. The
intermediate arrays therefore introduce no further multiplicities.

\paragraph{Independent bits and tensor entries.}
Transfer the exterior bits in the order specified by the splits and
concatenations. The source and destination share a processor header
for all pieces sent along one contact. Their array positions and bit
offsets have length $(\log h)^{O(1)}$. When a sphere collects bits
at its exterior station, the station is frozen during the ensuing
steps. When the blocks are reordered, permute its lists with them.
Neighborhood collection transports descriptions. Each of the $ed$
exterior bits travels only along its specified route, at most once
per transfer. The $O(\log h)$ steps therefore use
$ed(\log h)^{O(1)}$ payload work.

At the terminal stage, guess one bit per internal wire at its
designated endpoint and send it to the other endpoint in contact
and bundle order. A table call consists of the station description
and its ordered port bits. The terminal payload costs
$L_j(\log h)^{O(1)}$ by Equation~\eqref{eq:round-terminal}.
A scalar withdrawn earlier issues its call at withdrawal, with just
the short description. The input scalars are fixed. If their product
is $c$, one nullary table of value $c^d$ supplies their contribution
to the entire expansion. Naming this element of $K$ costs $O(\log h)$
bits. For fixed summand choices, the free
payload guesses now range exactly over the terminal wire
assignments. The product of call values is their contraction
weight. Summing over these assignments and then over summands gives
the input tensor power on its ordered $ed$ exterior bits.

\paragraph{Size of the fixed network.}
At each stage Equation~\eqref{eq:round-terminal} gives
$HN_b\le mS h^{O_D(a)+o(1)}$. There are $O(\log h)$ steps, and
$N_b\ge R$. The bounds for processors and for their total charge are
therefore
\begin{align}
 P&\le mS h^{O_D(a)+o(1)}/R,\nonumber\\
 V&\le P h^{C_2}+mS h^{O_D(a)+o(1)}
     +(L_j+ed)h^{o(1)}+2^{(\log h)^{O(1)}}.\label{eq:asm-volume}
\end{align}
The first two terms pay for the common work per position and the
short records per occurrence. The third term pays for independent
wire bits, and the last for the root computations. Every local
record has charge at most $2^{(\log h)^{O(1)}}$, since $N_b\le S$,
$\log S=O((\log h)^2)$ and an occurrence has only polynomially many
bits in $h$. The exponent $C_2$ and the exponent bounding the
declaration count can be fixed independently of $M$ in
$R=h^{M+o(1)}$. Common work uses $O(h)$-bit addresses and
$r^{O_D(1)}$ candidate contacts. The dependence on $M$ is confined
to factors polynomial in $\log S$ and the additive root cost.

The longest dictionary call uses $O_D(T_j)+h^{O_D(a)+o(1)}$ bits.
Choose $a$ small enough that $Ca+2a<\beta$ for each coefficient $C$
in the short-work and label bounds, and $b+(1+\delta)a<1$ for all
allowed $b$. Equation~\eqref{eq:round-terminal} gives
$T_j=O(h^{b+(1+\delta/3)a})$, so the calls have length below $h$.
Use Lemma~\ref{lem:fixed-transcript} with $\tau=a\delta/3$, then
take $M$ sufficiently large to absorb both $P h^{C_0}$ and
$P h^{C_2+\tau}$. For the three other charges in
Equation~\eqref{eq:asm-volume}, the costs per $md$ before the
factor $h^\tau$ are
\[
 \frac{mS h^{Ca+o(1)}}{md}=h^{-b+Ca+o(1)},\qquad
 \frac{ed}{md}\le h^{-2a},\qquad
 \frac{L_j}{md}=O(h^{-a}\log h).
\]
Here $d=T_0S$ and $T_0=\Theta(h^b)$. The additive root cost is
$n^{o(1)}$ while $m\ge n$. Consequently the fixed network has size
at most $md\,h^{-a+\tau+o(1)}\le md\,h^{-(1-\delta)a}$ for
large $n$. Its maximum table arity, including the activity and
length fields, is at most $h^{b+(1+\delta)a}$. A dictionary table
retains its short length-class label, ready for another application
of Lemma~\ref{lem:macro}. This completes the proof of that lemma.
